% ============================================================
%  BÁO CÁO ĐỒ ÁN CSC10004
%  Restaurant Management System
%  Compile: pdflatex (2 lần) hoặc xelatex
% ============================================================
\documentclass[12pt, a4paper]{article}

% ── Tiếng Việt ──────────────────────────────────────────────
\usepackage[utf8]{inputenc}
\usepackage[T5]{fontenc}
\usepackage[vietnamese]{babel}

% ── Layout ──────────────────────────────────────────────────
\usepackage[top=2.5cm, bottom=2.5cm, left=3cm, right=2cm]{geometry}
\usepackage{setspace}
\onehalfspacing

% ── Bảng & hình ─────────────────────────────────────────────
\usepackage{booktabs}
\usepackage{tabularx}
\usepackage{longtable}
\usepackage{array}
\usepackage{multirow}
\usepackage{graphicx}
\usepackage{float}
\usepackage{xcolor}
\definecolor{codebg}{RGB}{248,248,248}
\definecolor{primary}{RGB}{0,70,127}
\definecolor{accent}{RGB}{160,25,25}
\definecolor{warn}{RGB}{200,100,0}

% ── Hộp chú thích ───────────────────────────────────────────
\usepackage{tcolorbox}
\tcbuselibrary{skins}
\newtcolorbox{notebox}{colback=warn!8, colframe=warn!60, fonttitle=\bfseries,
  title=Lưu ý kỹ thuật, left=4pt, right=4pt, top=3pt, bottom=3pt}
\newtcolorbox{bugbox}{colback=red!6, colframe=red!55, fonttitle=\bfseries,
  title=Bug xác nhận từ test, left=4pt, right=4pt, top=3pt, bottom=3pt}
\newtcolorbox{passbox}{colback=green!6, colframe=green!50!black, fonttitle=\bfseries,
  title=Kết quả test, left=4pt, right=4pt, top=3pt, bottom=3pt}

% ── Code listing ────────────────────────────────────────────
\usepackage{listings}
\lstset{
  backgroundcolor=\color{codebg},
  basicstyle=\ttfamily\footnotesize,
  keywordstyle=\color{primary}\bfseries,
  commentstyle=\color{gray!70}\itshape,
  stringstyle=\color{accent},
  breaklines=true,
  frame=single,
  numbers=left,
  numberstyle=\tiny\color{gray},
  language=C++,
  tabsize=4,
  showstringspaces=false,
  xleftmargin=1.5em
}

% ── TikZ sơ đồ ──────────────────────────────────────────────
\usepackage{tikz}
\usetikzlibrary{shapes.geometric, arrows.meta, positioning, fit, backgrounds, calc}

% ── Math ────────────────────────────────────────────────────
\usepackage{amsmath}
\usepackage{amssymb}

% ── Hyperlink ───────────────────────────────────────────────
\usepackage[hidelinks]{hyperref}

% ── Header / footer ─────────────────────────────────────────
\usepackage{fancyhdr}
\pagestyle{fancy}
\fancyhf{}
\lhead{\small CSC10004 -- Cấu trúc dữ liệu và Giải thuật}
\rhead{\small Restaurant Management System}
\cfoot{\thepage}

% ── Định dạng tiêu đề ────────────────────────────────────────
\usepackage{titlesec}
\titleformat{\section}{\large\bfseries\color{primary}}{{\thesection.}}{0.5em}{}[\titlerule]
\titleformat{\subsection}{\normalsize\bfseries\color{primary!80}}{{\thesubsection.}}{0.5em}{}
\titleformat{\subsubsection}{\normalsize\bfseries}{{\thesubsubsection.}}{0.5em}{}

% ────────────────────────────────────────────────────────────
\begin{document}

% ============================================================
%  TRANG BÌA
% ============================================================
\begin{titlepage}
  \centering
  \vspace*{0.5cm}
  {\large\bfseries ĐẠI HỌC QUỐC GIA TP. HỒ CHÍ MINH}\\[0.2cm]
  {\large\bfseries TRƯỜNG ĐẠI HỌC KHOA HỌC TỰ NHIÊN}\\[0.1cm]
  {\large KHOA CÔNG NGHỆ THÔNG TIN}\\[1.5cm]

  \rule{\textwidth}{1.5pt}\\[0.4cm]
  {\LARGE\bfseries BÁO CÁO ĐỒ ÁN MÔN HỌC}\\[0.3cm]
  {\Large\bfseries\color{primary} Restaurant Management System}\\[0.2cm]
  {\normalsize Ứng dụng Quản lý Nhà hàng sử dụng\\
  Cấu trúc dữ liệu và Giải thuật tự cài đặt}\\[0.3cm]
  \rule{\textwidth}{1.5pt}\\[1.5cm]

  \begin{minipage}[t]{0.48\textwidth}
    \raggedright
    \textbf{Môn học:}\\
    CSC10004 -- Cấu trúc dữ liệu và Giải thuật\\[0.5cm]
    \textbf{Lớp:} CQ2025/4\\[0.5cm]
    \textbf{Giảng viên hướng dẫn:}\\
    ThS. Bùi Duy Đăng\\
    ThS. Huỳnh Lâm Hải Đăng\\
    ThS. Nguyễn Thanh Tình
  \end{minipage}
  \hfill
  \begin{minipage}[t]{0.48\textwidth}
    \raggedright
    \textbf{Nhóm sinh viên:}\\[0.3cm]
    \begin{tabular}{@{}ll@{}}
      \toprule
      \textbf{MSSV} & \textbf{Họ và tên} \\
      \midrule
      24120472 & Trương Huệ Trí \\
      24120415 & Vũ Lê Bá Phúc \\
      \bottomrule
    \end{tabular}\\[0.5cm]
    \textbf{Phân công:}\\[0.2cm]
    \begin{tabular}{@{}ll@{}}
      24120472 & Cài đặt thư viện (\texttt{lib/}) \\
      24120415 & Xây dựng ứng dụng (\texttt{app/}) \\
    \end{tabular}
  \end{minipage}

  \vfill
  {\large TP. Hồ Chí Minh, tháng 6 năm 2025}
\end{titlepage}

% ============================================================
%  MỤC LỤC
% ============================================================
\tableofcontents
\newpage

% ============================================================
%  1. GIỚI THIỆU ĐỀ TÀI
% ============================================================
\section{Giới thiệu đề tài}

\subsection{Bối cảnh}

Đồ án môn học CSC10004 yêu cầu mỗi cặp sinh viên: (1) tự cài đặt thư viện C++ bao gồm các cấu
trúc dữ liệu và giải thuật đã học, và (2) xây dựng ứng dụng thực tế sử dụng ít nhất 5 trong 7
thành phần thư viện đó một cách có ý nghĩa, không hời hợt.

\subsection{Ứng dụng được xây dựng}

Nhóm xây dựng \textbf{Restaurant Management System} -- hệ thống quản lý nhà hàng chạy trên
giao diện dòng lệnh (console), viết bằng C++17. Hệ thống giải quyết bốn bài toán nghiệp vụ:

\begin{enumerate}
  \item Xác thực và quản lý tài khoản theo vai trò (Admin / Staff / Customer).
  \item Quản lý thực đơn: thêm, xóa, tìm kiếm và hiển thị theo nhiều tiêu chí.
  \item Quản lý đơn hàng: tạo, tra cứu, cập nhật trạng thái, thống kê doanh thu.
  \item Lập lịch xử lý tác vụ bếp theo chính sách ưu tiên đa tầng.
\end{enumerate}

Mỗi bài toán sử dụng một cấu trúc dữ liệu hoặc giải thuật cụ thể từ thư viện tự cài đặt, phù hợp
với từng chức năng -- minh họa cách áp dụng DSA trong một ứng dụng thực tế.

% ============================================================
%  2. YÊU CẦU VÀ PHẠM VI
% ============================================================
\section{Yêu cầu và phạm vi hệ thống}

\subsection{Vai trò người dùng}

\begin{itemize}
  \item \textbf{Admin:} Quản lý tài khoản, quản lý thực đơn, xem toàn bộ đơn hàng và thống kê.
  \item \textbf{Staff:} Tiếp nhận và cập nhật trạng thái đơn hàng, điều phối hàng chờ bếp.
  \item \textbf{Customer:} Xem thực đơn, đặt món, theo dõi trạng thái đơn của mình.
\end{itemize}

\subsection{Phạm vi kỹ thuật}

\begin{itemize}
  \item Giao diện: console -- không yêu cầu GUI theo đề bài.
  \item Ngôn ngữ: C++17. Các CTDL chính tự cài đặt trong \texttt{lib/}.
        \texttt{std::vector} chỉ dùng làm bộ nhớ nội bộ bên trong thư viện và làm kiểu trả
        về tạm thời -- không phải CTDL chính mà người dùng thao tác trực tiếp.
  \item Lưu trữ: tệp văn bản thuần (\texttt{users.txt}, \texttt{menu.txt}, \texttt{orders.txt}),
        nạp khi khởi động và lưu sau các thao tác.
  \item Tổ chức mã: thư viện trong \texttt{lib/}, ứng dụng trong \texttt{app/}, tách biệt
        theo yêu cầu đề bài.
\end{itemize}

\subsection{Cấu trúc thư mục}

\begin{lstlisting}[language=bash, numbers=none, caption=Tổ chức thư mục mã nguồn]
DSA_Application/
  lib/
    LinkedList.hpp   Stack.hpp        Queue.hpp
    PriorityQueue.hpp  BST.hpp        AVL.hpp
    HashTable.hpp    Algorithms.hpp   Detail.hpp
  app/
    models/          # User, MenuItem, Order, OrderItem, KitchenTask
    managers/        # UserManager, MenuManager, OrderManager, KitchenManager
    ui/              # ConsoleUI
    utils/           # FileIO, StringUtils, TimeUtils, AlgorithmHelper
  data/
    users.txt  menu.txt  orders.txt
  main.cpp
\end{lstlisting}

% ============================================================
%  3. ÁNH XẠ NGHIỆP VỤ SANG DSA
% ============================================================
\section{Ánh xạ nghiệp vụ sang cấu trúc dữ liệu và giải thuật}

\subsection{Bảng tổng hợp}

\begin{table}[H]
\centering
\renewcommand{\arraystretch}{1.5}
\begin{tabularx}{\textwidth}{|p{2.8cm}|p{3.1cm}|p{2.8cm}|X|}
\hline
\textbf{Chức năng} & \textbf{Vấn đề dữ liệu} & \textbf{DSA chọn} & \textbf{Lý do chọn$^*$} \\
\hline
Đăng nhập / xác thực & Tra cứu tài khoản theo username & \texttt{HashTable} (AVL chaining) & Tra cứu theo khoá duy nhất -- O(1) trung bình theo lý thuyết \\
\hline
Thêm/xóa/sửa tài khoản & CRUD trên tập tài khoản & \texttt{HashTable} & insert / find / erase O(1) trung bình theo lý thuyết \\
\hline
Load thực đơn từ file & Sắp xếp toàn bộ danh sách & \texttt{QuickSort} & O($n \log n$) trung bình theo lý thuyết \\
\hline
Thêm món mới vào menu & Chèn vào danh sách gần đã có thứ tự & \texttt{InsertionSort} & O($n$) trên dữ liệu gần sắp xếp theo lý thuyết \\
\hline
Sắp xếp menu theo giá & Sắp xếp khi hiển thị & \texttt{HeapSort} & O($n \log n$) theo lý thuyết; dùng trong AlgorithmHelper \\
\hline
Sắp xếp đơn theo thời gian & Lịch sử đơn mới nhất trước & \texttt{MergeSort} & O($n \log n$) theo lý thuyết; stable sort \\
\hline
Tìm món theo tên & Tìm trong mảng đã sắp xếp & Binary Search (thủ công) & O($\log n$) theo lý thuyết trên mảng sắp xếp \\
\hline
Tạo / tìm đơn hàng & Tra cứu theo \texttt{orderId} & \texttt{BST} & O($\log n$) trung bình, O($n$) xấu nhất (xem Mục~3.3) \\
\hline
Xem lịch sử đơn hàng & Xuất theo thứ tự \texttt{orderId} tăng dần & BST inorder traversal & Tính chất BST: inorder cho khoá tăng dần \\
\hline
Hàng chờ bếp & Chọn đơn theo ưu tiên đa tầng & \texttt{PriorityQueue} (heap) & top() O(1); push/pop O($\log n$) theo lý thuyết \\
\hline
\end{tabularx}
{\footnotesize $^*$ Tất cả độ phức tạp trong bảng là kết quả phân tích lý thuyết theo \cite{clrs}. Nhóm không thực hiện đo đạc hiệu năng thực nghiệm.}
\caption{Ánh xạ nghiệp vụ -- DSA trong Restaurant Management System}
\end{table}

\subsection{Phân tích chi tiết từng ánh xạ}

\subsubsection{HashTable cho quản lý tài khoản}

Chức năng đăng nhập yêu cầu tra cứu tài khoản theo \texttt{username}. Vì thao tác chủ yếu là
lookup một khoá cụ thể (không cần duyệt có thứ tự), \texttt{HashTable} phù hợp hơn mảng tuyến
tính: theo phân tích lý thuyết \cite{clrs}, hash table với separate chaining có độ phức tạp tra cứu
trung bình O(1) khi hệ số tải được kiểm soát.

Thư viện cài đặt \texttt{HashTable<Key,T>} theo phương pháp \emph{separate chaining}, trong đó
mỗi bucket là một cây \texttt{AVL}. Nội bộ, các bucket được lưu trong \texttt{std::vector}
(chi tiết cài đặt bộ nhớ, không phải CTDL chính mà người dùng thao tác).

\textbf{Các phương thức chính gồm:}
\texttt{insert}, \texttt{find}, \texttt{contains}, \texttt{erase}, \texttt{clear},
\texttt{operator[]}, \texttt{size}, \texttt{bucket\_count}.

\textbf{Ánh xạ nghiệp vụ:}
\begin{itemize}
  \item Đăng nhập $\rightarrow$ \texttt{find(username)} + kiểm tra mật khẩu.
  \item Tạo tài khoản $\rightarrow$ \texttt{insert(username, user)}.
  \item Xóa tài khoản $\rightarrow$ \texttt{erase(username)}.
  \item Đổi mật khẩu $\rightarrow$ \texttt{find(username)} rồi cập nhật trực tiếp trên pointer.
\end{itemize}

\begin{notebox}
  \texttt{HashTable} không hỗ trợ rehashing -- không có phương thức \texttt{rehash}.
  Nếu số lượng tài khoản tăng lớn hơn nhiều so với \texttt{BucketCount} (mặc định 16), hiệu
  suất có thể giảm do va chạm tích lũy. Ở quy mô demo, điều này không làm sai kết quả.
\end{notebox}

\subsubsection{Mảng động + Giải thuật sắp xếp cho thực đơn}

\texttt{MenuManager} lưu thực đơn trong mảng động tự quản lý (\texttt{MenuItem* items}, cấp phát
thủ công với \texttt{new[]} và \texttt{resize()}). Mảng được chọn vì thực đơn thường duyệt toàn
bộ và ít thay đổi, cung cấp truy cập O(1) theo chỉ mục và duyệt tuyến tính O($n$) hiệu quả.

\textbf{Chiến lược sắp xếp hai pha:}

\begin{table}[H]
\centering
\begin{tabularx}{\textwidth}{llX}
\toprule
\textbf{Kịch bản} & \textbf{Giải thuật} & \textbf{Hàm gọi} \\
\midrule
Load $n$ món từ file & \texttt{QuickSort} & \texttt{ds::quickSort(items, items+size, cmp)} \\
Thêm $k$ món mới ($k \ll n$) & \texttt{InsertionSort} & \texttt{ds::insertionSort(items, items+size, cmp)} \\
\bottomrule
\end{tabularx}
\end{table}

\textbf{Pivot của QuickSort} (\texttt{Algorithms.hpp}, dòng 101):

\begin{lstlisting}
auto pivot = *(first + (last - first) / 2);  // pivot la phan tu o giua
\end{lstlisting}

Pivot là \textbf{phần tử ở giữa} (middle element), không phải pivot ngẫu nhiên.

\textbf{AlgorithmHelper.hpp} -- ngoài hai giải thuật trên, ứng dụng còn sử dụng thêm:
\begin{itemize}
  \item \texttt{ds::heapSort}: sắp xếp menu theo giá giảm dần (cao $\to$ thấp).
  \item \texttt{ds::mergeSort}: sắp xếp đơn hàng theo thời gian (mới $\to$ cũ).
  \item \texttt{ds::selectionSort}: lấy top-$N$ món đắt nhất.
  \item \texttt{ds::bubbleSort}: sắp xếp danh sách nhỏ theo tên.
\end{itemize}

Ngoài ra, \texttt{ConsoleUI} gọi trực tiếp \texttt{ds::insertionSort} khi hiển thị danh sách món
lúc đặt hàng (\texttt{ConsoleUI.hpp}~dòng~714).

\textbf{Tìm kiếm theo tên:} \texttt{findByName} dùng binary search thủ công riêng
(không gọi \texttt{ds::binarySearch}), kỳ vọng O($\log n$) sau khi danh sách đã sắp xếp.

\subsubsection{BST cho quản lý đơn hàng}

Mỗi đơn hàng có \texttt{orderId} là khoá nguyên, duy nhất. Hệ thống cần: (1) tra cứu nhanh theo
ID, (2) xuất lịch sử theo thứ tự ID tăng dần, (3) tính doanh thu trên toàn bộ đơn.

BST cung cấp cả ba nhu cầu này qua cùng một cấu trúc. Theo phân tích lý thuyết \cite{clrs}:
\begin{itemize}
  \item \texttt{insert(order)}: O($\log n$) trung bình, O($n$) xấu nhất (cây lệch hoàn toàn).
  \item \texttt{find(orderId)}: O($\log n$) trung bình, O($n$) xấu nhất.
  \item \texttt{inorder(callback)}: O($n$) -- tính chất BST: inorder cho khoá theo thứ tự tăng dần.
  \item \texttt{remove(orderId)}: O($\log n$) trung bình.
\end{itemize}

\textbf{Các phương thức chính gồm:}
\texttt{insert}, \texttt{remove}, \texttt{find}, \texttt{contains}, \texttt{inorder},
\texttt{min}, \texttt{max}, \texttt{clear}, \texttt{size}, \texttt{empty}.
BST \textbf{không có} \texttt{preorder} hay \texttt{postorder} -- chỉ có \texttt{inorder}.

\begin{notebox}
  App tạo \texttt{orderId} tăng đơn điệu bắt đầu từ 1001. Với BST không tự cân bằng, chèn
  liên tiếp các khoá tăng dần sẽ tạo cây lệch hoàn toàn về phải, suy thoái thành O($n$) cho
  mọi thao tác. Ở quy mô demo (số đơn nhỏ), điều này không làm sai kết quả trong các test
  hiện có, nhưng là giới hạn kỹ thuật cần nhận thức.
  Hướng cải thiện: thay \texttt{BST} bằng \texttt{AVL} cho \texttt{OrderManager}.
\end{notebox}

\subsubsection{PriorityQueue cho hàng chờ bếp}

Bếp không xử lý đơn theo FIFO mà theo chính sách ưu tiên đa tầng. \texttt{PriorityQueue}
được dùng để quản lý hàng chờ này: luôn trả về task có độ ưu tiên cao nhất ở \texttt{top()}.

\textbf{Cài đặt thư viện:} \texttt{PriorityQueue<T, Compare>} dùng \texttt{std::vector<T>}
làm bộ nhớ heap nội bộ. Hai thao tác cốt lõi: \texttt{sift\_up} (sau push) và \texttt{sift\_down}
(sau pop).

\textbf{Các phương thức chính gồm:} \texttt{push}, \texttt{pop}, \texttt{top}, \texttt{empty},
\texttt{size}, \texttt{clear}.

\texttt{KitchenTask::operator>} định nghĩa thứ tự ưu tiên ba tầng:
\begin{lstlisting}
bool operator>(const KitchenTask& other) const {
    if (isVIP != other.isVIP)       return isVIP > other.isVIP;
    if (priority != other.priority) return priority < other.priority;
    return orderTime < other.orderTime;
}
\end{lstlisting}

\textbf{Comparator sau khi sửa} (\texttt{KitchenManager.hpp}):
\begin{lstlisting}
struct KitchenTaskComparator {
    // compare(a,b)=true nghia la a kem uu tien hon b => a bi day xuong
    // => dung b > a de task uu tien cao hon (b) noi len dung
    bool operator()(const KitchenTask& a, const KitchenTask& b) const {
        return b > a;
    }
};
\end{lstlisting}

\begin{notebox}
  Trong \texttt{ds::PriorityQueue}, semantics: \texttt{compare(a,b) = true} có nghĩa a sẽ bị
  đẩy xuống (a kém ưu tiên hơn b). Phiên bản ban đầu dùng \texttt{return a > b} đã tạo heap
  ngược (task ưu tiên thấp nổi lên đầu). Lỗi được phát hiện qua test và đã sửa thành
  \texttt{return b > a}. Kết quả sau khi sửa: 3/3 test ưu tiên PASS (xem Mục~\ref{sec:pq}).
\end{notebox}

% ============================================================
%  4. THIẾT KẾ HỆ THỐNG
% ============================================================
\section{Thiết kế hệ thống}

\subsection{Sơ đồ phụ thuộc module}

\begin{figure}[H]
\centering
\resizebox{0.96\textwidth}{!}{%
\begin{tikzpicture}[
  node distance=1.2cm and 1.4cm,
  box/.style={rectangle, draw, rounded corners=2pt, minimum width=2.4cm,
              minimum height=0.8cm, align=center, font=\small},
  lib/.style={box, fill=blue!8, draw=blue!45},
  mgr/.style={box, fill=green!8, draw=green!55!black},
  ui/.style={box,  fill=orange!10, draw=orange!65},
  file/.style={box, fill=gray!12, draw=gray!55},
  arr/.style={-Stealth, thick, gray!70}
]
\node[lib] (ht)   {HashTable};
\node[lib, right=1.6cm of ht]  (bst) {BST};
\node[lib, right=1.6cm of bst] (pq)  {PriorityQueue};
\node[lib, right=1.6cm of pq]  (alg) {Algorithms};
\node[lib, below=0.5cm of ht, xshift=0.8cm] (avl) {AVL};

\node[mgr, below=2cm of ht]           (um) {UserManager};
\node[mgr, right=1.6cm of um]         (mm) {MenuManager};
\node[mgr, right=1.6cm of mm]         (om) {OrderManager};
\node[mgr, right=1.6cm of om]         (km) {KitchenManager};

\node[ui,  below=1.2cm of mm, xshift=1.2cm] (cui) {ConsoleUI};
\node[file, below=1.2cm of cui] (fio) {FileIO};

\draw[arr] (avl)  -- (ht);
\draw[arr] (ht)   -- (um);
\draw[arr] (bst)  -- (om);
\draw[arr] (pq)   -- (km);
\draw[arr] (alg)  -- (mm);
\draw[arr] (alg)  to[bend right=15] (cui);

\draw[arr] (um) -- (cui);
\draw[arr] (mm) -- (cui);
\draw[arr] (om) -- (cui);
\draw[arr] (km) -- (cui);
\draw[arr] (cui) -- (fio);

\node[font=\footnotesize\bfseries, above=0.25cm of ht, xshift=3.2cm] {\texttt{lib/}};
\node[font=\footnotesize\bfseries, left=0.15cm of um] {\texttt{app/}};
\end{tikzpicture}}
\caption{Sơ đồ phụ thuộc module chính trong dự án}
\end{figure}

% ============================================================
%  5. SƠ ĐỒ MINH HỌA CÁC CTDL CHÍNH
% ============================================================
\section{Minh họa cấu trúc dữ liệu}

\subsection{HashTable với AVL separate chaining}

\begin{figure}[H]
\centering
\begin{tikzpicture}[
  bucket/.style={draw, minimum width=2.4cm, minimum height=0.6cm, font=\small},
  node/.style={draw, circle, minimum size=0.7cm, font=\footnotesize}
]
\node[bucket, fill=blue!8] (b0) at (0,0)   {bucket[0]};
\node[bucket, fill=blue!8] (b1) at (0,-0.8) {bucket[1]};
\node[bucket, fill=blue!8] (b2) at (0,-1.6) {bucket[2]};
\node[font=\small]          at (0,-2.2) {$\vdots$};
\node[bucket, fill=blue!8] (bn) at (0,-2.8) {bucket[15]};

\node[node, fill=green!15] (n0a) at (3.0, 0)    {\texttt{u1}};
\node[node, fill=green!15] (n1a) at (3.0, -0.8)  {\texttt{u2}};
\node[node, fill=green!15] (n1b) at (4.0, -0.8)  {\texttt{u3}};
\node[node, fill=green!15] (n1c) at (3.5, -1.5)  {\texttt{u5}};

\draw[-Stealth] (b0) -- (n0a) node[midway, above, font=\footnotesize] {AVL};
\draw[-Stealth] (b1) -- (n1a) node[midway, above, font=\footnotesize] {AVL};
\draw[-] (n1a) -- (n1b);
\draw[-] (n1a) -- (n1c);

\node[bucket, fill=gray!8] at (3.0,-2.8) {\small (empty)};
\draw[-Stealth] (bn) -- (3.0,-2.8);

\node[font=\footnotesize\itshape, below=0.1cm of bn, xshift=2.4cm]
  {Mỗi bucket là một cây AVL$<$pair$<$Key,T$>>$};
\end{tikzpicture}
\caption{HashTable với AVL separate chaining (16 bucket mặc định, không rehashing)}
\end{figure}

\subsection{BST lưu Order theo orderId -- inorder traversal}

\begin{figure}[H]
\centering
\begin{tikzpicture}[
  bnode/.style={draw, circle, minimum size=0.9cm, font=\small, fill=orange!12}
]
\node[bnode] (r) at (0,0) {1003};
\node[bnode] (l) at (-1.8,-1.2) {1001};
\node[bnode] (rr) at (1.8,-1.2) {1005};
\node[bnode] (lr) at (-0.9,-2.4) {1002};
\node[bnode] (rl) at (0.9,-2.4) {1004};

\draw (r)--(l); \draw (r)--(rr); \draw (l)--(lr); \draw (rr)--(rl);

\node[font=\footnotesize, right=0.2cm of rr] {$\leftarrow$ root nếu chèn theo thứ tự này};

\draw[-Stealth, dashed, gray!70] (-3.2,-3.2) -- (-0.4,-3.2)
  node[midway, below, font=\footnotesize] {inorder: 1001 1002 1003 1004 1005};
\end{tikzpicture}
\caption{BST minh họa inorder traversal -- trường hợp dữ liệu không tuần tự. \\
  \small Lưu ý: trong app thực, orderId được tạo tăng dần (1001, 1002, \ldots) nên BST có
  thể lệch phải hoàn toàn.}
\end{figure}

\subsection{Max-heap cho KitchenTask}

\begin{figure}[H]
\centering
\begin{tikzpicture}[
  hnode/.style={draw, rectangle, rounded corners=2pt, minimum width=3.6cm,
                minimum height=0.75cm, font=\footnotesize, fill=red!8}
]
\node[hnode] (top) at (0,0)        {VIP, pri=0, T=9:00};
\node[hnode] (l)   at (-3.2,-1.6)  {VIP, pri=1, T=9:05};
\node[hnode] (rr)  at (3.2,-1.6)   {non-VIP, pri=0, T=9:00};
\node[hnode] (ll)  at (-5.6,-3.4)  {non-VIP, pri=1, T=9:10};
\node[hnode] (lr)  at (-0.8,-3.4)  {VIP, pri=1, T=9:15};

\draw (top)--(l); \draw (top)--(rr); \draw (l)--(ll); \draw (l)--(lr);

\node[font=\footnotesize\itshape, below=0.4cm of ll, xshift=2.4cm]
  {Thứ tự pop lý thuyết: VIP, pri thấp, sớm nhất trước};
\end{tikzpicture}
\caption{Cấu trúc heap cho KitchenTask theo thứ tự ưu tiên dự kiến \\
  \small (thứ tự thực tế phụ thuộc vào chiều comparator -- xem Mục~\ref{sec:pq})}
\end{figure}

% ============================================================
%  6. PHÂN TÍCH TỪNG CTDL / GIẢI THUẬT
% ============================================================
\section{Phân tích từng cấu trúc dữ liệu và giải thuật}
\label{sec:analysis}

\subsection{HashTable với AVL separate chaining}

\paragraph{Độ phức tạp lý thuyết.}
\begin{table}[H]
\centering
\begin{tabular}{lccc}
\toprule
\textbf{Thao tác} & \textbf{Trung bình} & \textbf{Xấu nhất} & \textbf{Phương thức} \\
\midrule
\texttt{insert} & O(1) & O($\log n$)$^*$ & \texttt{HashTable.hpp:37} \\
\texttt{find}   & O(1) & O($\log n$)$^*$ & \texttt{HashTable.hpp:61} \\
\texttt{erase}  & O(1) & O($\log n$)$^*$ & \texttt{HashTable.hpp:75} \\
\bottomrule
\end{tabular}
\end{table}
{\footnotesize $^*$ Xấu nhất trong một bucket (AVL) khi tất cả khoá hash về cùng bucket.}

\subsection{QuickSort và InsertionSort}

\paragraph{QuickSort} (\texttt{Algorithms.hpp:96--125}): chia-và-trị, pivot là phần tử ở giữa
(\texttt{*(first + (last-first)/2)}). Theo \cite{clrs}: O($n \log n$) trung bình, O($n^2$) xấu nhất
(khi pivot luôn chia không đều -- ví dụ dữ liệu đã sắp xếp với pivot đầu/cuối; pivot giữa giảm
thiểu rủi ro này).

\paragraph{InsertionSort} (\texttt{Algorithms.hpp:42--55}): duyệt từ đầu đến cuối, chèn từng
phần tử vào đúng vị trí trong phần đã sắp xếp. Theo \cite{clrs}: O($n$) trên dữ liệu gần có thứ tự
(mỗi phần tử chỉ dịch chuyển ít bước), O($n^2$) xấu nhất.

\paragraph{Lý do chiến lược hai pha.}
QuickSort khi load $n$ món từ file (O($n \log n$) trung bình theo lý thuyết); InsertionSort khi
thêm $k$ món mới vào danh sách gần sắp xếp (O($n$) trung bình theo lý thuyết khi $k \ll n$).
Đây là lập luận dựa trên phân tích lý thuyết; nhóm không thực hiện đo đạc thực nghiệm.

\subsection{BST}

\paragraph{Độ phức tạp lý thuyết \cite{clrs}.}
\begin{table}[H]
\centering
\begin{tabular}{lccc}
\toprule
\textbf{Thao tác} & \textbf{Trung bình} & \textbf{Xấu nhất} & \textbf{Phương thức} \\
\midrule
\texttt{insert} & O($\log n$) & O($n$) & \texttt{BST.hpp:119} \\
\texttt{find}   & O($\log n$) & O($n$) & \texttt{BST.hpp:140} \\
\texttt{remove} & O($\log n$) & O($n$) & \texttt{BST.hpp:133} \\
\texttt{inorder} & O($n$) & O($n$) & \texttt{BST.hpp:174} \\
\bottomrule
\end{tabular}
\end{table}

Trường hợp xấu O($n$) xảy ra khi cây lệch hoàn toàn (ví dụ: chèn khoá tăng đơn điệu) --
đây là rủi ro được xác nhận qua test với orderId 1001..1007 (xem Mục~7.2).

\subsection{PriorityQueue}

\paragraph{Độ phức tạp lý thuyết.}
\begin{table}[H]
\centering
\begin{tabular}{lcc}
\toprule
\textbf{Thao tác} & \textbf{Độ phức tạp} & \textbf{Phương thức} \\
\midrule
\texttt{push} & O($\log n$) -- sift\_up & \texttt{PriorityQueue.hpp:62} \\
\texttt{top}  & O(1) & \texttt{PriorityQueue.hpp:60} \\
\texttt{pop}  & O($\log n$) -- sift\_down & \texttt{PriorityQueue.hpp:72} \\
\bottomrule
\end{tabular}
\end{table}

\paragraph{Ghi chú về comparator.}
Với \texttt{std::less<T>} (mặc định), \texttt{ds::PriorityQueue} là max-heap theo \cite{clrs}.
\texttt{KitchenTaskComparator} dùng \texttt{return b > a}: khi b ưu tiên hơn a, comparator trả
\texttt{true} → a bị đẩy xuống → b nổi lên đúng vị trí. Hành vi này đã được xác nhận qua test
suite: 3/3 test kiểm tra thứ tự ưu tiên ba tầng đều PASS (xem Mục~\ref{sec:pq}).

\subsection{AVL (dùng nội bộ trong HashTable)}

AVL (\texttt{lib/AVL.hpp}) là BST tự cân bằng: sau mỗi insert hoặc remove, cây thực hiện phép
xoay qua \texttt{rebalance()} để duy trì $|\text{balance\_factor}| \leq 1$. Theo \cite{clrs},
điều kiện này đảm bảo chiều cao O($\log n$) -- mọi thao tác find/insert/remove đều
O($\log n$) trong mọi trường hợp. Dùng làm bucket trong HashTable thay vì danh sách liên kết.

API: \texttt{insert}, \texttt{remove}, \texttt{find}, \texttt{contains}, \texttt{inorder},
\texttt{min}, \texttt{max}, \texttt{clear}, \texttt{size}.

% ============================================================
%  7. KIỂM THỬ
% ============================================================
\section{Kiểm thử thư viện}

Toàn bộ test được thực hiện qua file \texttt{test\_lib.cpp} -- gọi trực tiếp các hàm trong
\texttt{lib/} với input và expected output xác định, không qua UI.

\textbf{Compile và chạy:}
\begin{lstlisting}[language=bash, numbers=none]
g++ -std=c++17 -I. test_lib.cpp -o test_lib && ./test_lib
\end{lstlisting}

\textbf{Nguồn gốc test case.} Các test được xây dựng từ ba nguồn: (1) đặc tả API của mỗi cấu
trúc dữ liệu (các phương thức \texttt{insert}, \texttt{find}, \texttt{remove}, \texttt{size},
\texttt{clear}, v.v.\ được khai báo trong \texttt{lib/}); (2) tính chất cốt lõi theo lý thuyết
DSA (ví dụ: inorder traversal của BST trả khoá theo thứ tự tăng dần \cite{clrs}, PriorityQueue
luôn trả phần tử ưu tiên cao nhất); (3) các trường hợp biên phổ biến (phần tử không tồn tại,
cấu trúc rỗng, thao tác trùng lặp).

\textbf{Giới hạn của PASS.} Kết quả PASS xác nhận tính đúng đắn chức năng cho các trường hợp
đã kiểm tra. PASS không phản ánh hiệu năng, độ phức tạp thực nghiệm, hay hành vi với dữ liệu
ngoài phạm vi test. Mọi nhận định về Big-O trong báo cáo là phân tích lý thuyết có nguồn \cite{clrs}.

\begin{passbox}
\textbf{Kết quả tổng hợp: 70 PASS \ | \ 0 FAIL}\\
HashTable: 19/19 \textbf{PASS} \quad
BST: 17/17 \textbf{PASS} \quad
PriorityQueue: 11/11 \textbf{PASS} \quad
Algorithms: 16/16 \textbf{PASS} \quad
AVL: 9/9 \textbf{PASS}
\end{passbox}

\begin{figure}[H]
\centering
\includegraphics[width=0.88\textwidth]{img/00_summary.png}
\caption{Tổng hợp kết quả test suite -- chạy thực tế từ \texttt{test\_lib.exe}}
\end{figure}

\subsection{HashTable -- 19/19 PASS}

\begin{figure}[H]
\centering
\includegraphics[width=0.92\textwidth]{img/01_hashtable.png}
\caption{HashTable: toàn bộ 19 test PASS}
\end{figure}

\begin{longtable}{|p{4cm}|p{4.5cm}|c|}
\hline
\textbf{Test case} & \textbf{Kết quả mong đợi} & \textbf{Thực tế} \\
\hline
\texttt{size()} sau 3 inserts & 3 & \textbf{PASS} \\
\hline
\texttt{find("alice")} & != nullptr, value == 10 & \textbf{PASS} \\
\hline
\texttt{find("nobody")} & nullptr & \textbf{PASS} \\
\hline
\texttt{contains("bob")} / \texttt{contains("ghost")} & true / false & \textbf{PASS} \\
\hline
Duplicate insert \texttt{"alice"} với value mới & size không đổi, value cập nhật & \textbf{PASS} \\
\hline
\texttt{erase("bob")} & true, size giảm, find trả nullptr & \textbf{PASS} \\
\hline
\texttt{erase("ghost")} (không tồn tại) & false, size không đổi & \textbf{PASS} \\
\hline
Update password qua pointer & \texttt{*find() = newpass} hoạt động & \textbf{PASS} \\
\hline
\texttt{find} user không tồn tại (login guard) & nullptr & \textbf{PASS} \\
\hline
\texttt{clear()} & size=0, find trả nullptr & \textbf{PASS} \\
\hline
\end{longtable}

\subsection{BST -- 17/17 PASS}

\begin{figure}[H]
\centering
\includegraphics[width=0.92\textwidth]{img/02_bst.png}
\caption{BST: 17/17 PASS -- bao gồm cả demo degeneration với orderId tuần tự}
\end{figure}

\begin{longtable}{|p{4.2cm}|p{5.2cm}|c|}
\hline
\textbf{Test case} & \textbf{Kết quả mong đợi} & \textbf{Thực tế} \\
\hline
Insert 5 đơn (không tuần tự), \texttt{find(1003)} & value == 300.0 & \textbf{PASS} \\
\hline
\texttt{find(9999)} & nullptr & \textbf{PASS} \\
\hline
Inorder (5 khóa không tuần tự) & 1001 1002 1003 1004 1005 (tăng dần) & \textbf{PASS} \\
\hline
Tính tổng qua inorder lambda & 1500.0 & \textbf{PASS} \\
\hline
Remove node lá (1001) & true, find=nullptr, size=4 & \textbf{PASS} \\
\hline
Remove node một con (1002) & true, find=nullptr, size=3 & \textbf{PASS} \\
\hline
Remove node hai con (1003, gốc) & true, successor thay thế đúng & \textbf{PASS} \\
\hline
Inorder sau 3 remove & \texttt{\{1004,1005\}} & \textbf{PASS} \\
\hline
\texttt{remove(9999)} & false, size không đổi & \textbf{PASS} \\
\hline
Insert tuần tự 1001..1007, inorder & \texttt{\{1001..1007\}} đúng & \textbf{PASS} \\
\hline
\end{longtable}

\begin{notebox}
\textbf{Ghi nhận degeneration (từ test output):}\\
Chèn tuần tự orderId 1001..1007 tạo cây lệch hoàn toàn về phải. Chiều cao = 7 = O($n$).
Inorder vẫn đúng về mặt logic, nhưng \texttt{find()} phải duyệt 7 bước thay vì $\lceil \log_2 7 \rceil = 3$.
\end{notebox}

\subsection{PriorityQueue -- 11/11 PASS}
\label{sec:pq}

\subsubsection{Sanity check: max-heap với \texttt{std::less<int>} -- 4/4 PASS}

\begin{figure}[H]
\centering
\includegraphics[width=0.75\textwidth]{img/03a_pq_sanity.png}
\caption{PriorityQueue với \texttt{std::less<int>}: max-heap hoạt động đúng}
\end{figure}

\begin{lstlisting}[numbers=none, caption=Output thực tế từ test\_lib.exe]
push {3,1,4,1,5} -> top() == 5       [PASS]
pop order = {5,4,3,1,1}              [PASS]  (max-heap xac nhan)
pop() on empty -> still empty        [PASS]
\end{lstlisting}

Lý do test: kiểm tra tính chất heap cơ bản (phần tử lớn nhất luôn ở đầu) theo định nghĩa
trong \cite{clrs}.

\subsubsection{KitchenTask -- bug phát hiện, sửa, xác nhận: 3/3 PASS}

\textbf{Bug phát hiện.} Phiên bản ban đầu của \texttt{KitchenTaskComparator} dùng
\texttt{return a > b}. Trong \texttt{ds::PriorityQueue}, \texttt{compare(a,b)=true} có nghĩa
a bị đẩy xuống. Vì \texttt{a > b} là true khi a ưu tiên hơn b, task ưu tiên cao bị đẩy
xuống -- tạo heap ngược. Quan sát từ test (xem ảnh minh họa):

\begin{figure}[H]
\centering
\includegraphics[width=0.85\textwidth]{img/03b_pq_bug.png}
\caption{Minh họa hành vi sai khi dùng \texttt{return a > b}: non-VIP nổi lên đầu heap}
\end{figure}

\textbf{Sửa và xác nhận.} Đổi comparator thành \texttt{return b > a}
(\texttt{KitchenManager.hpp}, đã áp dụng vào production code).

\begin{figure}[H]
\centering
\includegraphics[width=0.92\textwidth]{img/03c_pq_fixed.png}
\caption{Sau khi sửa (\texttt{return b > a}): thứ tự ưu tiên ba tầng đúng, 3/3 PASS}
\end{figure}

\begin{lstlisting}[numbers=none, caption=Output thực tế sau khi sửa comparator]
[PASS] VIP on top over non-VIP    top(): orderId=1003  isVIP=1
[PASS] VIP pri=0 on top over pri=1  top(): orderId=2001  priority=0
[PASS] Earlier orderTime on top    top(): orderId=3001  orderTime=1000

Full pop order:
#1003(VIP=1,pri=0,t=1000) #1004(VIP=1,pri=0,t=1010)
#1005(VIP=1,pri=1,t=1000) #1001(VIP=0,pri=0,t=1000)
#1002(VIP=0,pri=0,t=1005)
\end{lstlisting}

Thứ tự pop khớp với yêu cầu nghiệp vụ: VIP trước, giai đoạn thấp hơn trước, đặt sớm hơn trước.

\subsection{Algorithms -- 16/16 PASS}

\begin{figure}[H]
\centering
\includegraphics[width=0.85\textwidth]{img/04_algorithms.png}
\caption{6 giải thuật sắp xếp + 2 tìm kiếm -- 16/16 PASS}
\end{figure}

\begin{longtable}{|p{4.5cm}|p{4cm}|c|}
\hline
\textbf{Test case} & \textbf{Kết quả} & \textbf{Thực tế} \\
\hline
\texttt{quickSort} \{5,2,8,1,9,3,7,4,6,0\} & 0 1 2 3 4 5 6 7 8 9 & \textbf{PASS} \\
\hline
\texttt{insertionSort} (same input) & 0 1 2 3 4 5 6 7 8 9 & \textbf{PASS} \\
\hline
\texttt{heapSort} (same input) & 0 1 2 3 4 5 6 7 8 9 & \textbf{PASS} \\
\hline
\texttt{mergeSort} (same input) & 0 1 2 3 4 5 6 7 8 9 & \textbf{PASS} \\
\hline
\texttt{selectionSort} (same input) & 0 1 2 3 4 5 6 7 8 9 & \textbf{PASS} \\
\hline
\texttt{bubbleSort} (same input) & 0 1 2 3 4 5 6 7 8 9 & \textbf{PASS} \\
\hline
\texttt{quickSort} với cmp ngược (desc) & 9 8 7 6 5 4 3 2 1 0 & \textbf{PASS} \\
\hline
\texttt{insertionSort} dữ liệu gần sắp xếp & đúng thứ tự & \textbf{PASS} \\
\hline
\texttt{quickSort} chuỗi ký tự & ``Banh Mi'' ``Bun Bo'' ``Com Tam'' \ldots & \textbf{PASS} \\
\hline
\texttt{linearSearch} phần tử tồn tại & tìm thấy 7 & \textbf{PASS} \\
\hline
\texttt{linearSearch} phần tử không có & trả \texttt{end} & \textbf{PASS} \\
\hline
\texttt{binarySearch} phần tử tồn tại (7, 1, 15) & tìm thấy đúng & \textbf{PASS} \\
\hline
\texttt{binarySearch} phần tử không có (6) & trả \texttt{last} gốc & \textbf{PASS} \\
\hline
\end{longtable}

\begin{notebox}
\textbf{Sửa đã áp dụng cho \texttt{ds::binarySearch}:}
Phiên bản ban đầu trả về \emph{insertion point} (vị trí hẹp dần của \texttt{last}) khi không
tìm thấy, thay vì \texttt{last} gốc. Đã sửa bằng cách lưu \texttt{original\_last} trước vòng
lặp và trả về nó khi không tìm thấy. Sau sửa: PASS.\\
Lưu ý: \texttt{MenuManager::findByName} dùng binary search thủ công riêng (trả \texttt{nullptr}
khi không tìm thấy), không gọi \texttt{ds::binarySearch} -- sửa này không ảnh hưởng đến app.
\end{notebox}

\subsection{AVL -- 9/9 PASS}

\begin{figure}[H]
\centering
\includegraphics[width=0.85\textwidth]{img/05_avl.png}
\caption{AVL: insert tuần tự vẫn cân bằng, remove đúng -- 9/9 PASS}
\end{figure}

\begin{longtable}{|p{4.5cm}|p{4cm}|c|}
\hline
\textbf{Test case} & \textbf{Kết quả} & \textbf{Thực tế} \\
\hline
Insert tuần tự 1..7, inorder & \{1,2,3,4,5,6,7\} & \textbf{PASS} \\
\hline
\texttt{find(4)}, \texttt{find(99)} & 4 / nullptr & \textbf{PASS} \\
\hline
\texttt{remove(4)}, inorder sau đó & \{1,2,3,5,6,7\} & \textbf{PASS} \\
\hline
\texttt{remove(99)} & false & \textbf{PASS} \\
\hline
\end{longtable}

% ============================================================
%  8. ĐÁNH GIÁ MỨC ĐỘ ĐÁP ỨNG
% ============================================================
\section{Đánh giá mức độ đáp ứng yêu cầu}

\subsection{Thành phần thư viện được sử dụng}

\begin{table}[H]
\centering
\renewcommand{\arraystretch}{1.3}
\begin{tabularx}{\textwidth}{clp{3.4cm}X}
\toprule
\textbf{\#} & \textbf{Thành phần} & \textbf{Nơi sử dụng} & \textbf{Vai trò} \\
\midrule
1 & \texttt{HashTable<K,V>} & \texttt{UserManager} & Đăng nhập, CRUD tài khoản \\
2 & \texttt{AVL<T>} & Nội bộ \texttt{HashTable} & Separate chaining per bucket \\
3 & \texttt{Algorithms.hpp} & \texttt{MenuManager}, \texttt{AlgorithmHelper}, \texttt{ConsoleUI} &
  QuickSort, InsertionSort, HeapSort, MergeSort, SelectionSort, BubbleSort \\
4 & \texttt{BST<T>} & \texttt{OrderManager} & CRUD đơn hàng + inorder traversal \\
5 & \texttt{PriorityQueue<T>} & \texttt{KitchenManager} & Lập lịch ưu tiên hàng chờ bếp \\
\midrule
-- & \texttt{LinkedList<T>} & -- & Không sử dụng \\
-- & \texttt{Stack<T>} & -- & Không sử dụng \\
-- & \texttt{Queue<T>} & -- & Không sử dụng \\
\bottomrule
\end{tabularx}
\caption{5/7 thành phần thư viện được sử dụng trong các module ứng dụng}
\end{table}

\subsection{Bug phát hiện qua test và đã sửa}

\begin{enumerate}
  \item \textbf{KitchenTaskComparator bị ngược chiều.}
        Phiên bản ban đầu (\texttt{return a > b}) tạo heap ngược: non-VIP nổi lên đầu.
        Phát hiện qua test: khi push VIP và non-VIP vào cùng queue, \texttt{top()} trả non-VIP.
        \textbf{Đã sửa:} đổi thành \texttt{return b > a} trong \texttt{KitchenManager.hpp}.
        \textbf{Xác nhận:} 3/3 test ưu tiên PASS sau khi sửa (xem Mục~\ref{sec:pq}).

  \item \textbf{\texttt{ds::binarySearch} trả insertion point khi không tìm thấy.}
        Phiên bản ban đầu hẹp dần biến \texttt{last} trong vòng lặp và trả giá trị đã hẹp,
        không phải \texttt{last} gốc. Phát hiện qua test: tìm 6 trong \{1,3,5,7,...\}
        trả vị trí của 7, không phải \texttt{end}.
        \textbf{Đã sửa:} lưu \texttt{original\_last} trước vòng lặp, trả về nó khi không tìm thấy
        (\texttt{lib/Algorithms.hpp}).
        \textbf{Xác nhận:} PASS sau khi sửa. Không ảnh hưởng đến app vì
        \texttt{MenuManager} dùng binary search thủ công riêng.
\end{enumerate}

\subsection{Hạn chế kỹ thuật còn tồn tại}

\begin{enumerate}
  \item \textbf{BST suy thoái với orderId tăng đơn điệu.}
        Quan sát từ test: chèn 1001..1007 tuần tự tạo cây lệch hoàn toàn về phải -- mọi
        thao tác find phải duyệt tất cả 7 nút. Theo lý thuyết \cite{clrs}, trường hợp này
        có chiều cao O($n$). Inorder vẫn cho kết quả đúng. Giải pháp: thay BST bằng AVL
        trong \texttt{OrderManager}.

  \item \textbf{HashTable không có rehashing.}
        Số bucket cố định (mặc định 16). Ở quy mô demo, số tài khoản nhỏ nên không làm sai
        kết quả trong các test hiện có.

  \item \textbf{Mọi nhận định về Big-O là lý thuyết, không phải đo đạc thực nghiệm.}
        Nhóm không thực hiện benchmark hiệu năng. Tất cả độ phức tạp được trích dẫn từ
        \cite{clrs} và áp dụng theo lý thuyết cấu trúc code.
\end{enumerate}

% ============================================================
%  KẾT LUẬN
% ============================================================
\section*{Kết luận}
\addcontentsline{toc}{section}{Kết luận}

Restaurant Management System minh họa cách áp dụng DSA trong một ứng dụng thực tế: mỗi bài
toán nghiệp vụ sử dụng một cấu trúc dữ liệu hoặc giải thuật phù hợp với từng chức năng.
Hệ thống sử dụng 5/7 thành phần thư viện trong các module nghiệp vụ: \texttt{UserManager}
dùng \texttt{HashTable}, \texttt{OrderManager} dùng \texttt{BST}, \texttt{KitchenManager}
dùng \texttt{PriorityQueue}, \texttt{MenuManager}/\texttt{AlgorithmHelper}/\texttt{ConsoleUI}
dùng các thuật toán trong \texttt{Algorithms.hpp}, và \texttt{HashTable} dùng \texttt{AVL}
cho các bucket.

Quá trình kiểm thử phát hiện hai bug (comparator PriorityQueue bị ngược chiều;
\texttt{ds::binarySearch} trả sai giá trị khi không tìm thấy) -- cả hai đã được sửa và xác nhận
lại qua test suite (70/70 PASS). Rủi ro kỹ thuật còn lại (BST suy thoái với orderId tăng đơn
điệu) được ghi nhận là hạn chế thiết kế, không làm sai kết quả trong các test hiện có.

Mọi nhận định về độ phức tạp trong báo cáo là phân tích lý thuyết theo \cite{clrs}; nhóm không
thực hiện đo đạc hiệu năng thực nghiệm.

% ============================================================
%  TÀI LIỆU THAM KHẢO
% ============================================================
\begin{thebibliography}{9}
\bibitem{clrs}
  T.~H. Cormen, C.~E. Leiserson, R.~L. Rivest, C.~Stein.
  \textit{Introduction to Algorithms}, 4th ed.
  MIT Press, 2022.

\bibitem{stroustrup}
  B.~Stroustrup.
  \textit{The C++ Programming Language}, 4th ed.
  Addison-Wesley, 2013.

\bibitem{cppreference}
  cppreference.com -- C++ Templates.
  \url{https://en.cppreference.com/w/cpp/language/templates}
\end{thebibliography}

\end{document}
